\documentclass[a4paper,10pt]{article}
\usepackage{cmap}
\usepackage[T2A]{fontenc}
\usepackage[utf8]{inputenc}
\usepackage[russian]{babel}
\usepackage{amsmath,amssymb,microtype}
\usepackage[a4paper,landscape,top=19mm,bottom=16mm,left=15mm,right=15mm,headheight=14pt,headsep=5mm]{geometry}
\usepackage{enumitem,fancyhdr,multicol}
\usepackage[unicode,hidelinks]{hyperref}
\hypersetup{pdftitle={Демоверсия к КР1 по теории вероятностей},pdfauthor={hse-se-lecture-notes}}
\pagestyle{fancy}
\fancyhf{}
\fancyhead[L]{\small\itshape ПИ, ФКН НИУ ВШЭ}
\fancyhead[C]{\small\itshape hse-se-lecture-notes}
\fancyhead[R]{\small\itshape Теория вероятностей, 2026/27}
\fancyfoot[C]{\small\thepage}
\renewcommand{\headrulewidth}{0pt}
\setlength{\parindent}{0pt}
\setlength{\columnsep}{10mm}
\raggedcolumns
\setlength{\abovedisplayskip}{7pt}
\setlength{\belowdisplayskip}{7pt}
\setlist[enumerate,1]{label=\arabic*.,leftmargin=6mm,labelsep=2mm,itemsep=10pt,parsep=0pt,topsep=0pt}
\newcommand{\PP}{\mathsf P}
\newcommand{\EE}{\mathsf E}
\newcommand{\DD}{\mathsf D}
\newcommand{\R}{\mathbb R}
\newcommand{\normalise}{Найти постоянную $C$ и функцию распределения $F_\xi(x)$.}
\begin{document}
\begin{multicols}{2}
\begin{center}
{\large\mdseries Демоверсия к КР1 по теории вероятностей}
\end{center}
\vspace{4pt}
\begin{enumerate}
% Блок 1: классическая и геометрическая вероятность.
\item\begin{minipage}[t]{\linewidth} В лотерее 15 билетов, из которых 5 выигрышные. Наугад выбирают три билета без возвращения. Найти вероятность того, что ровно два из них выигрышные.
\end{minipage}

\item\begin{minipage}[t]{\linewidth} В ящике 6 стандартных и 4 нестандартные детали. Наугад берут три детали без возвращения. Найти вероятность того, что хотя бы одна из выбранных деталей нестандартная.
\end{minipage}

\item\begin{minipage}[t]{\linewidth} В круге радиуса 2 равномерно выбирают случайную точку. Найти вероятность того, что расстояние от этой точки до центра круга больше $1{,}5$.
\end{minipage}

\item\begin{minipage}[t]{\linewidth} Два человека договорились встретиться между 12:00 и 14:00. Каждый приходит в случайный момент, равномерно распределённый на этом промежутке, ждёт 30 минут и уходит. Времена прихода независимы. Найти вероятность встречи.

% Блок 2: вариант 5, задача 1; вариант 6, задача 1.
\end{minipage}

\item\begin{minipage}[t]{\linewidth} Вероятность того, что случайно выбранный пирожок с вишней, равна $\PP(A)=3/8$, а вероятность того, что он сладкий, равна $\PP(B)=1/7$. События $A$ и $B$ независимы. Найти вероятность того, что пирожок с вишней или сладкий, и вероятность того, что пирожок с вишней, если известно, что он сладкий.
\end{minipage}

\item\begin{minipage}[t]{\linewidth} События $A$ и $B$ независимы, причём $\PP(A)=2/9$, $\PP(B)=3/10$. Найти $\PP(A\cup B)$ и $\PP(A\mid B)$.

% Блок 2: вариант 5, задача 2; вариант 6, задача 2.
\end{minipage}

\item\begin{minipage}[t]{\linewidth} Стрелок стреляет по мишени до первого попадания. Вероятность попадания при одном выстреле равна $p=2/7$, результаты выстрелов независимы. Найти вероятность того, что потребуется ровно 3 выстрела.
\end{minipage}

\item\begin{minipage}[t]{\linewidth} Монету бросают до первого выпадения орла. Вероятность выпадения орла при каждом броске равна $p=3/8$, броски независимы. Найти вероятность того, что потребуется ровно 2 броска.

% Блок 2: вариант 5, задачи 3--4.
\end{minipage}

\item\begin{minipage}[t]{\linewidth}\label{libraries} Три библиотеки хранят художественную и научную литературу. Доля научных книг в первой библиотеке равна $q_1=1/5$, во второй --- $q_2=1/10$, в третьей --- $q_3=1/20$. Библиотека выбирается с вероятностями
\[
p_1=\frac12,\qquad p_2=\frac14,\qquad p_3=\frac14.
\]
Из выбранной библиотеки наугад берут одну книгу. Найти вероятность того, что книга научная.
\end{minipage}

\item\begin{minipage}[t]{\linewidth} В условиях задачи~\ref{libraries} найти вероятность того, что выбранная научная книга взята из первой библиотеки.

% Блок 2: вариант 6, задачи 3--4. Данные подтверждены пользователем.
\end{minipage}

\item\begin{minipage}[t]{\linewidth}\label{factories} Три завода поставляют детали. Доли заводов в поставке равны
\[
p_1=\frac13,\qquad p_2=\frac12,\qquad p_3=\frac16.
\]
Вероятности брака для деталей с этих заводов соответственно равны
\[
q_1=\frac1{15},\qquad q_2=\frac1{30},\qquad q_3=\frac1{10}.
\]
Найти вероятность того, что случайно выбранная из поставки деталь бракованная.
\end{minipage}

\item\begin{minipage}[t]{\linewidth} В условиях задачи~\ref{factories} найти вероятность того, что бракованная деталь произведена первым заводом.

% Блок 6: новая задача на телескопическую сумму.
\end{minipage}

\item\begin{minipage}[t]{\linewidth} Распределение дискретной случайной величины $\xi$ задано формулой
\[
\PP\{\xi=k\}=\frac{C}{k(k+2)},\qquad k=1,2,\ldots
\]
Найти: а) постоянную $C$; б) вероятность $\PP\{\xi\geq4\}$.

% Блок 3: плотности из вариантов 1, 2, 8, 7.

Решить аналогичную задачу для распределения
\[
\PP\{\xi=k\}=\frac{C}{(k+1)(k+4)},\quad k=0,1,\ldots,
\]
найдя $C$ и $\PP\{\xi\geq3\}$.
\end{minipage}

\item\begin{minipage}[t]{\linewidth} Плотность случайной величины $\xi$ задана формулой
\[
p_\xi(x)=\begin{cases}\dfrac{C}{(x+1)^2},&x\geq0,\\[3pt]0,&x<0.\end{cases}
\]
\normalise
\end{minipage}

\item\begin{minipage}[t]{\linewidth} Плотность случайной величины $\xi$ задана формулой
\[
p_\xi(x)=\begin{cases}C\cos x,&-\pi/2\leq x\leq\pi/2,\\0,&\text{иначе}.\end{cases}
\]
\normalise
\end{minipage}

\item\begin{minipage}[t]{\linewidth} Плотность случайной величины $\xi$ задана формулой
\[
p_\xi(x)=\begin{cases}Cx^2+\dfrac32x,&0\leq x\leq1,\\[3pt]0,&\text{иначе}.\end{cases}
\]
\normalise
\end{minipage}

\item\begin{minipage}[t]{\linewidth} Плотность случайной величины $\xi$ задана формулой
\[
p_\xi(x)=\begin{cases}C+3x^2-2x^3,&0\leq x\leq1,\\0,&\text{иначе}.\end{cases}
\]
\normalise

% Блок 4: новая задача с двумя пунктами.
\end{minipage}

\item\begin{minipage}[t]{\linewidth} На вероятностном пространстве $\Omega=[0,1]$ с $\sigma$-алгеброй измеримых по Лебегу множеств и мерой Лебега в качестве вероятности найти функцию распределения случайной величины $\xi$, если:
\[
\text{а) }\xi(\omega)=3\omega-1;
\]
\[
\text{б) }\xi(\omega)=\begin{cases}2\omega,&0\leq\omega\leq\tfrac12,\\3(1-\omega),&\tfrac12<\omega\leq1.\end{cases}
\]

% Блок 3: графики вариантов 1, 2, 8, 7; координаты исходных рисунков.
\end{minipage}
\item\begin{minipage}[t]{\linewidth} На пространстве $\Omega=[0,1]$ с мерой Лебега в качестве вероятности задана величина
\[
\xi(\omega)=\begin{cases}3\omega,&0\leq\omega\leq\tfrac13,\\2-3\omega,&\tfrac13<\omega\leq\tfrac23,\\6\omega-4,&\tfrac23<\omega\leq1.\end{cases}
\]
Найти функцию распределения и плотность $\xi$.
\end{minipage}


\item\begin{minipage}[t]{\linewidth} Случайная величина $\xi$ имеет плотность
\[
p_\xi(x)=\frac12e^{-|x|},\qquad x\in\R.
\]
Найти функцию распределения и плотность величины $\eta=\xi^2$.

% Блок 5, задача 1; блок 7, задача 2: дискретные моменты.
\end{minipage}

\item\begin{minipage}[t]{\linewidth} Случайная величина $\xi$ принимает значения $2k+1$ с вероятностями
\[
\PP\{\xi=2k+1\}=\frac34\left(\frac14\right)^k,
\qquad k=0,1,2,\ldots
\]
Найти математическое ожидание и дисперсию $\xi$.
\end{minipage}

\item\begin{minipage}[t]{\linewidth} Распределение случайной величины $\xi$ задано формулой
\[
\PP\{\xi=k\}=e^{-3}\frac{3^k}{k!},\qquad k=0,1,2,\ldots
\]
Найти математическое ожидание и дисперсию величины $\eta=2\xi-1$.
Моменты $\xi$ вычислить непосредственно через суммы.

% Блок 7, задача 3; блок 5, задачи 2--3; блок 7, задача 4.
\end{minipage}

\item\begin{minipage}[t]{\linewidth} Случайная величина $\xi$ имеет функцию распределения
\[
F_\xi(x)=\begin{cases}0,&x<1,\\1-e^{-2(x-1)},&x\geq1.\end{cases}
\]
Найти плотность, математическое ожидание и дисперсию $\xi$.
\end{minipage}

\item\begin{minipage}[t]{\linewidth} Плотность случайной величины $\xi$ задана формулой
\[
p_\xi(x)=\begin{cases}\dfrac{64}{x^5},&x\geq2,\\[3pt]0,&x<2.\end{cases}
\]
Найти математическое ожидание и дисперсию величин $\xi$ и $\eta=2/\xi$.
\end{minipage}

\item\begin{minipage}[t]{\linewidth} Случайная величина $\xi$ имеет плотность
\[
p_\xi(x)=\frac1{3\sqrt{2\pi}}\exp\left(-\frac{(x-2)^2}{18}\right),
\qquad x\in\R.
\]
Найти математическое ожидание и дисперсию $\xi$ и $\eta=4-2\xi$, используя переход к стандартному нормальному распределению и свойства линейных преобразований.
\end{minipage}

\item\begin{minipage}[t]{\linewidth} Плотность случайной величины $\xi$ задана формулой
\[
p_\xi(x)=\begin{cases}\dfrac2{(1+x)^3},&x\geq0,\\[3pt]0,&x<0.\end{cases}
\]
Вычислить $\EE\xi$ и $\EE\xi^2$, если они конечны. Определить, имеет ли $\xi$ конечную дисперсию.
\end{minipage}


\item\begin{minipage}[t]{\linewidth} Совместная плотность величин $\xi,\eta$ имеет вид
\[
p_{\xi,\eta}(u,v)=\begin{cases}C(u+2v),&0\leq u\leq2,\ 0\leq v\leq1,\\0,&\text{иначе}.\end{cases}
\]
Найти $C$, одномерные плотности $p_\xi$, $p_\eta$ и вероятность $\PP\{\xi+\eta\leq1\}$. Проверить независимость $\xi$ и $\eta$.
\end{minipage}


\item\begin{minipage}[t]{\linewidth} Точка $(\xi,\eta)$ равномерно распределена в треугольнике с вершинами $(0,0)$, $(3,0)$ и $(0,2)$. Найти совместную плотность, плотности обеих координат и вероятность $\PP\{\xi\leq1,\eta\leq1\}$.
\end{minipage}


\item\begin{minipage}[t]{\linewidth} Случайные величины $\xi,\eta$ независимы и имеют плотности
\[
p_\xi(x)=2e^{-2x},\qquad p_\eta(x)=3e^{-3x},\qquad x\geq0
\]
(при $x<0$ плотности равны нулю). Найти $\PP\{|\xi-\eta|\leq\tfrac12\}$.
\end{minipage}


\item\begin{minipage}[t]{\linewidth} Совместная плотность величин $\xi,\eta$ задана формулой
\[
p_{\xi,\eta}(u,v)=\begin{cases}\dfrac{192}{\pi(u^2+v^2)^4},&u^2+v^2\geq4,\\[3pt]0,&\text{иначе}.\end{cases}
\]
Найти функцию распределения и плотность величины $R=\sqrt{\xi^2+\eta^2}$, а также $\EE R$ и $\DD R$.
\end{minipage}


\item\begin{minipage}[t]{\linewidth} Случайные величины $\xi$ и $\eta$ независимы, их функции распределения равны $F$ и $G$. Выразить через $F$ и $G$ функции распределения $\max\{\xi,2\eta+1\}$ и $\min\{\xi,2\eta+1\}$.
\end{minipage}


\item\begin{minipage}[t]{\linewidth} Независимые величины $\xi,\eta$ имеют показательные распределения с параметрами $1$ и $2$ соответственно. Найти функции распределения и плотности величин $\max\{\xi,\eta^2\}$ и $\min\{\xi,\eta^2\}$.
\end{minipage}


\item\begin{minipage}[t]{\linewidth} Индикаторы $I_1,\ldots,I_n$ независимы, причём $\PP\{I_i=1\}=p$, $\PP\{I_i=0\}=1-p$. За успех начисляют 3 балла, за неудачу снимают 1 балл. Итоговый балл равен
\[
S=\sum_{i=1}^n(4I_i-1).
\]
Найти $\EE S$, $\DD S$ и $\EE S^2$, используя представление в виде суммы. Здесь $n\geq1$, $0<p<1$.
\end{minipage}


\item\begin{minipage}[t]{\linewidth} Величины $\xi,\eta$ независимы, причём
\[
\EE\xi=-1,\quad\EE\eta=3,\quad\DD\xi=2,\quad\DD\eta=5.
\]
Найти: а) $\EE(2\xi^2+\eta^2-3\xi\eta+\xi-2\eta)$;
б) $\EE(2\xi-\eta+1)^2$; в) $\DD(3\xi+2\eta-4)$.
\end{minipage}


\item\begin{minipage}[t]{\linewidth} Величины $\xi\in\{-2,0,2\}$ и $\eta\in\{-1,0,1\}$ удовлетворяют условиям
\[
\PP\{\xi\eta=0\}=1,\qquad\PP\{\xi=2\}=\PP\{\xi=-2\}=\frac16,
\]
\[
\PP\{\eta=1\}=\PP\{\eta=-1\}=\frac14.
\]
Восстановить таблицу совместного распределения. Найти ожидания, дисперсии, ковариацию и коэффициент корреляции. Проверить независимость и некоррелированность.
\end{minipage}


\item\begin{minipage}[t]{\linewidth} Правильную игральную кость бросают $n\geq1$ раз независимо. Пусть $\xi$ --- число выпадений чётных чисел, а $\eta$ --- число выпадений чисел, больших 3. Найти $\operatorname{cov}(\xi,\eta)$ и коэффициент корреляции. Независимы ли $\xi$ и $\eta$?
\end{minipage}


\item\begin{minipage}[t]{\linewidth} Случайные величины $\xi,\eta$ независимы,
\[
\PP\{\xi=2\}=\PP\{\xi=-2\}=\frac12,
\]
\[
\PP\{\eta=1\}=\PP\{\eta=-1\}=\frac13,\qquad\PP\{\eta=0\}=\frac13.
\]
Построить совместную таблицу распределения $\xi\eta$ и $\eta$. Будут ли эти величины независимыми? Некоррелированными?
\end{minipage}
\end{enumerate}
\end{multicols}
\end{document}
